\documentclass[11pt]{article}
\usepackage[letterpaper,margin=1in]{geometry}
\usepackage{amsmath,amssymb,amsthm,booktabs,microtype,xurl}
\usepackage[hidelinks]{hyperref}
\newtheorem{theorem}{Theorem}
\newtheorem{lemma}[theorem]{Lemma}
\newtheorem{corollary}[theorem]{Corollary}
\newtheorem{proposition}[theorem]{Proposition}
\newcommand{\C}{\mathbb C}
\newcommand{\Q}{\mathbb Q}
\newcommand{\R}{\mathbb R}
\newcommand{\ii}{\mathrm i}
\newcommand{\SIC}{\operatorname{SIC}}
\newcommand{\Hess}{\operatorname{Hess}}

\title{A 14-variable unipotent Keller map and\\
every-order image and vanishing obstructions}
\author{Alec Kriebel\\[2pt]\large with heavy assistance from ChatGPT 5.6 Sol}
\date{Research draft prepared 22 July 2026\\Not peer reviewed}

\begin{document}
\maketitle

\begin{abstract}
Starting from the recently announced three-variable counterexample to the
Jacobian conjecture, we give one self-contained consequence mechanism built
from a formal inverse pencil and a linear observable.  Its first realization
is an exact rational vector $g\in\Q[Z]^{14}$ with 24 monomials such that
$\det(I+sJg)=1$, $Jg$ is generically regular nilpotent, and $I+g$ has an exact
reduced fiber of three rational points.  The same inverse series gives an
explicit every-exponent failure of the Special Image Conjecture $\SIC(14)$.
After homogenization and symmetric reduction it gives a homogeneous
degree-eight Hessian-nilpotent polynomial in 30 variables for which
$\Delta^mP^{m+1}\ne0$ at every order.  A complementary rank-compressed cubic
route gives a homogeneous quartic in 44 variables with the same every-order
nonvanishing, directly addressing the classical quartic Vanishing
Conjecture.  The real 14-variable map also supplies direct counterexamples to
the Campbell, Chamberland--Meisters, and Kulikov injectivity and fixed-point
formulations.  These are one consequence cascade, not independent
discoveries.  All claims are accompanied by exact sparse certificates and
bounded-memory checkers.
\end{abstract}

\noindent\textbf{Verification disclaimer.}
The human author is a complete amateur exploring the limits of AI-assisted
mathematics and cannot independently verify these claims.  This is an
unreviewed artifact for expert checking, not an established result.  The
algebraic checks do not establish the literature or priority assertions.

\section{The unified result}

Let $E_n:\C[\xi,Z]\to\C[Z]$ be the linear map
\begin{equation}
 E_n(\xi^\alpha p(Z))=\partial_Z^\alpha p(Z),
 \label{eq:E}
\end{equation}
and let $\Delta_N$ denote the Laplacian in $N$ variables.  The paper proves
the following package directly; none of the operator statements is inferred
only from a logical equivalence with the Jacobian conjecture.

\begin{theorem}[Master theorem]\label{thm:master}
There are explicit exact objects with the following properties.
\begin{enumerate}
\item A vector $g\in\Q[Z]^{14}$ with 24 monomials satisfies
\[
 \det(I+sJg)=1,\qquad (Jg)^{14}=0,\qquad (Jg)^{13}\ne0.
\]
The map $T=I+g$ has generic degree three, geometric monodromy $S_3$, and an
exact scheme-theoretic fiber consisting of three reduced rational points.
\item For $A=-\xi\cdot g$ and $b=x+y+u_{11}$,
\[
 E_{14}(A^m)=0,\qquad E_{14}(bA^m)\ne0\qquad(m\ge1).
\]
Thus $\SIC(14)$ is false, with one fixed multiplier obstructing every
positive exponent.
\item There is a 28-variable nonhomogeneous Hessian-nilpotent polynomial
$P_{28}$ of degree eight and a 30-variable homogeneous Hessian-nilpotent
polynomial $P_{30}$ of degree eight such that, for $P=P_{28},P_{30}$,
\[
 \Delta^m P^m=0\quad(m\ge1),\qquad
 \Delta^mP^{m+1}\ne0\quad(m\ge0).
\]
The expanded term counts are 178 and 608, respectively.
\item There is a 44-variable homogeneous Hessian-nilpotent quartic $P_{44}$,
with 538 expanded monomials, such that
\[
 \Delta_{44}^mP_{44}^m=0\quad(m\ge1),\qquad
 \Delta_{44}^mP_{44}^{m+1}\ne0\quad(m\ge0).
\]
\end{enumerate}
\end{theorem}

The 30- and 44-variable endpoints are complementary: the former has smaller
ambient dimension, while the latter lies in the degree-four homogeneous case
used in Zhao's equivalence with the Jacobian conjecture~\cite{zhao-hnp}.  No global
minimality claim is made for either.

\section{The source map and one inverse series}

The following map is the announced three-dimensional counterexample that
initiated the present constructions~\cite{alpoge}.  Put $u=1+xy$ and define
\begin{align}
 F_1&=u^3z+y^2u(4+3xy),\nonumber\\
 F_2&=y+3xu^2z+3xy^2(4+3xy),\nonumber\\
 F_3&=2x-3x^2y-x^3z.\label{eq:F}
\end{align}
We use the identity-linear normalization
\begin{equation}
 \Phi=(F_3/2,F_2,F_1).
 \label{eq:Phi}
\end{equation}
Direct differentiation gives $\det J\Phi=1$.  The three points
\begin{equation}
 r_0=(0,0,-\tfrac14),\quad
 r_+=(1,-\tfrac32,\tfrac{13}2),\quad
 r_-=(-1,\tfrac32,\tfrac{13}2)
 \label{eq:sourcepoints}
\end{equation}
map to $(0,0,-1/4)$.

The common one-variable engine is the unique series $\tau\in t\Q[[t]]$
satisfying
\begin{equation}
 \tau+t\tau^3=\frac t2.
 \label{eq:tau}
\end{equation}
Let $D=1+3t\tau^2$.  Differentiating~\eqref{eq:tau} gives
$t\tau'=\tau/D$.  Set
\begin{equation}
 \bar x=t\tau'=\frac{\tau}{D},\qquad
 \bar y=-3t\tau,\qquad
 \bar z=tD(1-30t\tau^2-18t^2\tau^4).
 \label{eq:ray}
\end{equation}

\begin{lemma}[Exact inverse ray]\label{lem:ray}
The series in~\eqref{eq:ray} satisfy
\[
 \Phi(\bar x,\bar y,\bar z)=(t/2,0,t).
\]
\end{lemma}

\begin{proof}
Here $1+\bar x\bar y=1/D$.  Direct substitution into~\eqref{eq:F}
gives $F_1=t$, $F_2=0$, and
$F_3=2\tau(1+t\tau^2)=t$, the last equality being~\eqref{eq:tau}.
\end{proof}

Two linear observables below produce the series
\begin{align}
 q(t)&=\tau'-3\tau-\frac{\tau'-1/2}{t}
      =\sum_{m\ge0}q_mt^m,\label{eq:qdef}\\
 r(t)&=\tau'-3\tau+3t\tau\tau'
      =\sum_{m\ge0}r_mt^m.\label{eq:rdef}
\end{align}

\begin{lemma}[Coefficient-complete obstruction]\label{lem:coefficients}
For every $k\ge0$,
\begin{align}
q_{3k}=r_{3k}
 &=(-1)^k\frac{\binom{3k+1}{k}}{2^{2k+1}},\nonumber\\
q_{3k+1}=r_{3k+1}
 &=(-1)^{k+1}
 \frac{3\binom{3k+1}{k}}{(3k+1)2^{2k+1}},\nonumber\\
q_{3k+2}
 &=(-1)^k\frac{\binom{3k+4}{k+1}}{2^{2k+3}},\label{eq:qcoeff}\\
r_{3k+2}
 &=(-1)^k\frac{3\binom{3k+2}{k}}{2^{2k+2}}.\label{eq:rcoeff}
\end{align}
In particular, $q_mr_m\ne0$ for every $m\ge0$.
\end{lemma}

\begin{proof}
Lagrange inversion in~\eqref{eq:tau} gives
\[
 [t^{3k+1}]\tau=
 (-1)^k\frac{\binom{3k+1}{k}}{(3k+1)2^{2k+1}}.
\]
Differentiation gives the residue class $0$, multiplication by $-3$ gives
class $1$, and the remaining summands in~\eqref{eq:qdef} and
\eqref{eq:rdef} give class $2$.  These classes are disjoint; elementary
binomial simplification gives~\eqref{eq:qcoeff}--\eqref{eq:rcoeff}.
\end{proof}

\section{Two inverse-series transfer principles}

The next two identities are the conceptual spine of the paper.

\subsection{The image-operator transfer}

Let $H\in\C[Z]^n$, suppose $\det(I+tJH)=1$, and let $Q_t$ be the formal
inverse of $I+tH$.  Put
\begin{equation}
 A_H(\xi,Z)=-\sum_{j=1}^n\xi_jH_j(Z).
\end{equation}

\begin{lemma}[Abhyankar--Gurjar transfer]\label{lem:image-transfer}
For every polynomial $p$,
\begin{equation}
 p(Q_t(Z))=
 \sum_{m\ge0}\frac{t^m}{m!}E_n\bigl(pA_H^m\bigr)(Z).
 \label{eq:image-transfer}
\end{equation}
Consequently $E_n(A_H^m)=0$ for every $m\ge1$, and any nonzero coefficient
$[t^m]p(Q_t)(Z_*)$ proves $E_n(pA_H^m)(Z_*)\ne0$.
\end{lemma}

\begin{proof}
Apply the Abhyankar--Gurjar inversion formula~\cite{bcw,zhao-ag} to
$I+tH=I-(-tH)$.
The determinant factor is one, and the multinomial theorem collects terms of
total $\xi$-degree $m$.  Taking $p=1$ gives the first conclusion.
\end{proof}

\subsection{The symmetric Laplacian transfer}

Suppose additionally that $H(0)=0$.  For $A,B\in\C^n$, define the
de Bondt--van den Essen symmetric transform~\cite{debondt}
\begin{equation}
 P_H(A,B)=\ii\sum_{j=1}^n H_j(A+\ii B)B_j,
 \qquad \Gamma_t=I-t\nabla P_H.
 \label{eq:symmetrization}
\end{equation}

\begin{lemma}[Symmetric transfer]\label{lem:symmetric-transfer}
Write $X=A+\ii B$.  In coordinates $(X,B)$,
\begin{equation}
 \Gamma_t(X,B)=
 \bigl(X+tH(X),(I+tJH(X)^T)B-\ii tH(X)\bigr).
 \label{eq:triangular}
\end{equation}
Hence
\[
 \det J\Gamma_t=\det(I+tJH)^2=1,
\]
so $P_H$ is Hessian nilpotent.  Moreover, the $A+\ii B$ component of
$\Gamma_t^{-1}$ is $(I+tH)^{-1}$.

Let $\lambda$ be a linear form on $\C^n$, put
$\ell(A,B)=\lambda(A+\ii B)$, let $v$ be its coefficient vector in the
$2n$ variables, and write $D_v=v\mathbin{\cdot}\nabla$.  If
\[
 \lambda((I+tH)^{-1}(Y_*))=\sum_{k\ge0}c_kt^k,
\]
then
\begin{equation}
 D_v\bigl(\Delta_{2n}^mP_H^{m+1}\bigr)(Y_*,0)
 =2^m m!(m+1)!c_{m+1}.
 \label{eq:laplace-transfer}
\end{equation}
\end{lemma}

\begin{proof}
Differentiating~\eqref{eq:symmetrization} gives
\[
 \nabla_A P_H=\ii JH(X)^TB,\qquad
 \nabla_B P_H=-JH(X)^TB+\ii H(X),
\]
which proves~\eqref{eq:triangular}.  The determinant identity for every
$t$ is equivalent to nilpotence of $\Hess P_H$.

Because $H(0)=0$, the polynomial $P_H$ has order at least two.  Zhao's
inversion formula~\cite[Theorem~3.4]{zhao-hnp} for a Hessian-nilpotent
polynomial therefore writes
\begin{equation}
 \Gamma_t^{-1}(Z)=Z+t\nabla R_t(Z),\qquad
 R_t=\sum_{m\ge0}
 \frac{t^m}{2^m m!(m+1)!}\Delta_{2n}^mP_H^{m+1}.
 \label{eq:zhao}
\end{equation}
Apply $\ell$ at $(Y_*,0)$ and compare the coefficient of $t^{m+1}$.
\end{proof}

Zhao also proved that Hessian nilpotence is equivalent to
$\Delta^mP^m=0$ for every $m\ge1$.  Thus a coefficient-complete inverse
observable gives a coefficient-complete counterexample to the corresponding
Vanishing Conjecture without ever expanding a high Laplacian
power~\cite{zhao-hnp,zhao-general-vc}.

\section{The 14-variable regular-nilpotent map}

Write
\begin{equation}
 \Phi=X+H_2+H_3+\cdots+H_7.
\end{equation}
The nonzero homogeneous pieces of $\Phi-I$ are
\begin{center}
\begin{tabular}{c|ccc}
\toprule
$d$&$H_{d,1}$&$H_{d,2}$&$H_{d,3}$\\
\midrule
2&0&$3xz$&$4y^2$\\
3&$-\tfrac32x^2y$&$12xy^2$&$3xyz$\\
4&$-\tfrac12x^3z$&$6x^2yz$&$7xy^3$\\
5&0&$9x^2y^3$&$3x^2y^2z$\\
6&0&$3x^3y^2z$&$3x^2y^4$\\
7&0&0&$x^3y^3z$\\
\bottomrule
\end{tabular}
\end{center}

Take state coordinates
\[
 U=(u_{11},u_{12};u_{21},u_{22},u_{23},u_{24};
 u_{31},u_{32},u_{33},u_{34},u_{35}).
\]
Let $B\in\operatorname{Mat}_{3\times11}(\Q)$ select the three chain heads,
so $BU=(u_{11},u_{21},u_{31})^T$, let
$N=J_2\oplus J_4\oplus J_5$ with ones immediately above the diagonal, and
put
\[
C=(H_{3,1},H_{4,1};H_{3,2},H_{4,2},H_{5,2},H_{6,2};
H_{3,3},H_{4,3},H_{5,3},H_{6,3},H_{7,3})^T.
\]
Define
\begin{equation}
 g(X,U)=\bigl(H_2(X)+BU,-C(X)-NU\bigr).
 \label{eq:blockg}
\end{equation}
In the coordinate order $Z=(x,y,z,U)$, this is
\begin{align}
g_1&=u_{11},&
g_2&=3xz+u_{21},&
g_3&=4y^2+u_{31},\nonumber\\
g_4&=\tfrac32x^2y-u_{12},&
g_5&=\tfrac12x^3z,\nonumber\\
g_6&=-12xy^2-u_{22},&
g_7&=-6x^2yz-u_{23},\nonumber\\
g_8&=-9x^2y^3-u_{24},&
g_9&=-3x^3y^2z,\nonumber\\
g_{10}&=-3xyz-u_{32},&
g_{11}&=-7xy^3-u_{33},\nonumber\\
g_{12}&=-3x^2y^2z-u_{34},&
g_{13}&=-3x^2y^4-u_{35},\nonumber\\
g_{14}&=-x^3y^3z.\label{eq:g}
\end{align}

\begin{theorem}\label{thm:g14}
The vector~\eqref{eq:g} has 24 monomials and
\[
 \det(I+sJg)=1.
\]
Its generic nilpotency index is 14.  The fiber of $T=I+g$ over
$c=(0,0,-1/4,0,\ldots,0)$ consists scheme-theoretically of exactly three
reduced rational points, obtained by lifting~\eqref{eq:sourcepoints}.
The generic degree is three and the geometric monodromy is $S_3$.
\end{theorem}

\begin{proof}
The Jacobian pencil has block form
\[
 I+sJg=
 \begin{pmatrix}
 I_3+sJH_2&sB\\
 -sJC&I_{11}-sN
 \end{pmatrix}.
\]
Since $N^5=0$, a Schur complement gives
\begin{align*}
\det(I+sJg)
 &=\det\!\left(I_3+sJH_2+s^2B(I-sN)^{-1}JC\right)\\
 &=\det\!\left(I_3+\sum_{d=2}^7s^{d-1}JH_d(X)\right)
 =\det J\Phi(sX)=1.
\end{align*}
Cayley--Hamilton gives $(Jg)^{14}=0$, while exact multiplication gives
\[
 \bigl((Jg)^{13}\bigr)_{1,5}=-3x^6y^4z\ne0.
\]

On the stated target the state equations uniquely give
$U=(I-N)^{-1}C(X)$, and the first block becomes $\Phi(X)$.  The source-fiber
ideal has lexicographic Gr\"obner basis
\begin{equation}
 4z+1-27x^2,\qquad 2y+3x,\qquad x(x-1)(x+1).
 \label{eq:fiberbasis}
\end{equation}
Its quotient is reduced of length three, proving the exact fiber assertion.

More generally, over a target $(Y,V)$ the state block gives
\[
 U=(I-N)^{-1}(V+C(X)),\qquad
 \Phi(X)=Y-B(I-N)^{-1}V.
\]
Thus $T$ is stably equivalent to $\Phi\times I_{11}$ over the target space,
not merely on the displayed special fiber.

For completeness, a target $(a,b,c)$ of the unnormalized map $F$ has inverse
parameter $t$ satisfying
\[
 2at^3-bt^2+2t-c=0,
\]
with
\[
 y=-\frac{bt^2+3c-6t}{2t^2},\quad
 x=\frac{t}{1-ty},\quad
 z=\frac{2x-3x^2y-c}{x^3}.
\]
Thus the generic degree is at most three; the reduced three-point fiber and
\'{e}taleness give the reverse inequality.  The cubic discriminant is
\[
 -4(27a^2c^2-18abc+16a+b^3c-b^2),
\]
and the discriminant in $c$ of the parenthesized quadratic is
$-(12a-b^2)^3\ne0$.  It is therefore squarefree in $\C(a,b)[c]$ and
not a square in $\C(a,b,c)$.  Hence the transitive degree-three geometric
monodromy group is $S_3$.
\end{proof}

The weights
\begin{equation}
 \omega=(1,1,1;2,3;2,3,4,5;2,3,4,5,6)
 \label{eq:weights}
\end{equation}
satisfy $\operatorname{wt}(g_i)=\omega_i+1$.  If
$D_s=\operatorname{diag}(s^{\omega_i})$, then
\begin{equation}
 D_s^{-1}TD_s=I+sg\qquad(s\ne0).
 \label{eq:conjugacy}
\end{equation}
Thus the entire punctured pencil is linearly conjugate.  One diagonal
conjugate has integer coefficients; separately, the member $I+g/4$ has an
entirely integral three-point collision and common target $(0,0,-1,0,\ldots,0)$.
These are two different normal forms.

\section{Every-exponent SIC and the named consequence cascade}

Let
\begin{equation}
 A=-\sum_{j=1}^{14}\xi_jg_j,\qquad b=x+y+u_{11},
 \qquad Z_*=(\tfrac12,0,1,0,\ldots,0).
\end{equation}
For the inverse $Q_t$ of $I+tg$, state elimination and
Lemma~\ref{lem:ray} give
\begin{equation}
 Q_{t,x}(Z_*)=\tau',\quad
 Q_{t,y}(Z_*)=-3\tau,\quad
 Q_{t,u_{11}}(Z_*)=-\frac{\tau'-1/2}{t}.
 \label{eq:qcoordinates}
\end{equation}
The last equality is also the first target equation
$Q_x+tQ_{u_{11}}=1/2$.  Hence $b(Q_t)(Z_*)=q(t)$.

\begin{theorem}\label{thm:sic}
For every $m\ge1$,
\[
 E_{14}(A^m)=0,\qquad
 E_{14}(bA^m)(Z_*)=m!q_m\ne0.
\]
Consequently $\ker E_{14}$ is not a Mathieu subspace and $\SIC(14)$ is
false.
\end{theorem}

\begin{proof}
Apply Lemma~\ref{lem:image-transfer} with $H=g$ and use
Lemma~\ref{lem:coefficients}.
\end{proof}

Zhao proved~\cite{zhao-image}
\[
 \ker E_{14}=\sum_{j=1}^{14}(\xi_j-\partial_{Z_j})\C[\xi,Z].
\]
Following Liu--Sun's terminology~\cite{liusun}, the assertion that this image
is a Mathieu subspace is the Special Image Conjecture.  Thus
Theorem~\ref{thm:sic} also refutes Zhao's regular-sequence Image Conjecture,
of which SIC is this specialization.  It does not, by itself, reverse other
one-way implications in the literature.

The same map gives the following direct real and complex corollaries.  The
conjecture numbers in the first column are Campbell's~\cite{campbell}; the
other formulations are due to Chamberland--Meisters~\cite{chamberland} and
Kulikov~\cite{kulikov}.
\begin{center}
\begin{tabular}{p{0.31\textwidth}p{0.26\textwidth}p{0.32\textwidth}}
\toprule
Statement refuted&Map&Exact obstruction\\
\midrule
Campbell Conj.~3 (unipotent univalence)&$T$ on $\R^{14}$&$\operatorname{Spec}JT=\{1\}$ and three real preimages\\
Chamberland--Meisters Conj.~2.1&$T$ on $\R^{14}$&all eigenvalues have modulus one, hence a uniform gap from zero\\
Campbell Conj.~4 (stability / weak injectivity)&$-T$&spectrum constantly $\{-1\}$ and noninjectivity\\
Kulikov $JN(14)$ and Campbell Conj.~6&$N_c=c-g$&nilpotent derivative and exactly three fixed points\\
Campbell Conj.~5 (fixed-point form)&translate $N_c$ at $p_0$&the origin and two other fixed points; spectrum $\{0\}$\\
\bottomrule
\end{tabular}
\end{center}
These statements are one transformation family, not five independent
discoveries.  The classical dynamical Markus--Yamabe problems already had
older counterexamples; the present maps only give unusually rigid
constant-spectrum realizations.

Padding $g$ by zero coordinates proves the corresponding map, fixed-point,
and $\SIC(n)$ failures for every $n\ge14$.  The generic Jordan type becomes
$(14,1^{n-14})$.

\section{The 28- and 30-variable Hessian companions}

Directly apply~\eqref{eq:symmetrization} to $H=g$ and call the result
$P_{28}$.  It is nonhomogeneous of degrees two through eight and has 178
expanded monomials.  At $(Z_*,0)$, take
\[
 \ell_{28}(A,B)=b(A+\ii B).
\]
Lemma~\ref{lem:symmetric-transfer} and~\eqref{eq:qcoordinates} give
\begin{equation}
 D_{v_{28}}\bigl(\Delta_{28}^mP_{28}^{m+1}\bigr)(Z_*,0)
 =2^m m!(m+1)!q_{m+1}\ne0.
 \label{eq:P28}
\end{equation}

For the homogeneous version introduce $w$ and put
\begin{equation}
 h(Z,w)=\bigl(w^7g(Z/w),0\bigr)\in\Q[Z,w]^{15}.
 \label{eq:h15}
\end{equation}
It is homogeneous of degree seven and has 24 monomials.

\begin{proposition}\label{prop:hindex}
The matrix $Jh$ has exact generic nilpotency index 14 and Jordan type
$(14,1)$.
\end{proposition}

\begin{proof}
Write
\[
 Jh=\begin{pmatrix}W&q\\0&0\end{pmatrix},\qquad
 W=w^6Jg(Z/w),\qquad
 q=w^6\bigl(7g-Jg(Z/w)(Z/w)\bigr).
\]
Besides $(Jg)^{14}=0$, exact sparse multiplication gives
$(Jg)^{13}g=0$.  Hence $W^{14}=0$ and $W^{13}q=0$, so $(Jh)^{14}=0$.
On the other hand,
\[
 \bigl((Jh)^{13}\bigr)_{1,5}=-3w^{67}x^6y^4z\ne0.
\]
The only partition of 15 with largest part 14 is $(14,1)$.
\end{proof}

Let $P_{30}=P_h$ as in~\eqref{eq:symmetrization}.  This polynomial is
homogeneous of degree eight and has 608 expanded monomials.  Since the last
coordinate of $I+th$ is $w$, its inverse on the slice $w=1$ is exactly the
inverse of $I+tg$.  With target $(Z_*,1;0)$ and
$\ell_{30}(A,B)=b(A+\ii B)$, Lemma~\ref{lem:symmetric-transfer} gives
\begin{equation}
 D_{v_{30}}\bigl(\Delta_{30}^mP_{30}^{m+1}\bigr)(Z_*,1;0)
 =2^m m!(m+1)!q_{m+1}\ne0.
 \label{eq:P30}
\end{equation}
Equations~\eqref{eq:P28} and~\eqref{eq:P30}, together with Hessian
nilpotence, prove parts (3) of Theorem~\ref{thm:master}.

\section{The every-order homogeneous quartic}

We now retain degree four at the cost of a larger ambient dimension.  A
factor-reusing stable degree reduction in the classical
framework~\cite{bcw} gives
\begin{equation}
 \Psi=X+H_2+H_3:\C^{13}\to\C^{13},\qquad \det J\Psi=1,
 \label{eq:Psi13}
\end{equation}
in variables
\[
X=(x,y,z,g1a,g1b,g2a,g2b,g3a,g3b,g4a,g4b,g5a,g6b).
\]
For reproducibility, the six stable operations are summarized below.  A
two-factor step appends outputs $a+P,b+Q$ and subtracts the indicated
coefficient times their product from the target coordinate; one-factor steps
reuse the named existing output.
\begin{center}\small
\begin{tabular}{c|c|p{0.53\textwidth}|c}
\toprule
step&target&new or reused factors&coefficient\\
\midrule
1&$\Phi_1$&$x^2,\ xz+3y$; add $g1a,g1b$&$-1/2$\\
2&$\Phi_2$&$3x^2y,\ 2z+xyz+3y^2$; add $g2a,g2b$&$1$\\
3&$\Phi_2$&$xy,\ g2a\,z+3xg2b$; add $g3a,g3b$&$-1$\\
4&$\Phi_3$&$xy^2,\ 7y+3xz+3xy^2+x^2yz$; add $g4a,g4b$&$1$\\
5&$\Phi_3$&add $g5a+g4a\,xy$; reuse $g1b+xz+3y$&$-1$\\
6&$\Phi_3$&reuse $g3a+xy$; add $g6b+g4a\,g1b-yg4b$&$1$\\
\bottomrule
\end{tabular}
\end{center}
A final output shear subtracts the product of the $g3a$- and $g1b$-outputs
from the $g4b$-output.  Each displayed step is a composition with triangular
source and target automorphisms, so it preserves the constant Jacobian up to
its linear normalization.  Let $M$ be the resulting map and $L=JM(0)$, so
$\Psi=L^{-1}M$ and $\det J\Psi=1$.  Exact coefficient reduction factors
$H_3=BK$ with
$K=(K_1,\ldots,K_8)$ and
\begin{align}
K_1={}&\tfrac12x(g1a\,z+g1b\,x),\nonumber\\
K_2={}&g2a\,g3a\,z-3g2a\,y^2+3g2b\,g3a\,x+g3b\,xy+12xy^2,\nonumber\\
K_3={}&-g1b\,g3a\,g4a+g3a\,g4b\,y-3g4a\,xz+g5a\,xz-g6b\,xy+3xyz,\nonumber\\
K_4={}&3x^2y,\nonumber\\
K_5={}&2g1b\,g3a\,g4a-2g3a\,g4b\,y+6g4a\,xz-2g5a\,xz+2g6b\,xy-5xyz,
\label{eq:K}\\
K_6={}&xy^2,\nonumber\\
K_7={}&-g1b\,xy-7g2a\,g3a\,z+21g2a\,y^2-21g2b\,g3a\,x-g3a\,xz
-7g3b\,xy-84xy^2,\nonumber\\
K_8={}&g4a\,xy.\nonumber
\end{align}
Here
\begin{equation}
B(U_1,\ldots,U_8)=
(U_1,U_2,U_3,0,-3U_2,U_4,U_5,0,0,U_6,U_7,U_8,0).
\end{equation}
Define the cubic homogeneous vector in 22 variables
\begin{equation}
 h_4(X,U,w)=\bigl(wH_2(X)+w^2BU,-K(X),0\bigr).
 \label{eq:h22}
\end{equation}
A Schur complement gives
\[
 \det(I+sJh_4)
 =\det\!\left(I+s wJH_2+s^2w^2BJK\right)
 =\det J\Psi(swX)=1;
\]
hence $Jh_4$ is nilpotent.  The checker verifies this matrix identity
directly.

Set
\begin{equation}
Y=L^{-1}(\tfrac12,0,1,0,\ldots,0)
=(\tfrac12,0,1,0,0,0,-2,0,\ldots,0)\in\Q^{13},
\end{equation}
and $R_*=(Y,0^8,1)\in\Q^{22}$.  If $X(t)$ is the first block of
$(I+th_4)^{-1}(R_*)$, the second block gives $U=tK(X)$ and therefore
\[
 \Psi(tX)=tY,\qquad M(tX)=t(\tfrac12,0,1,0,\ldots,0).
\]
The stable outputs and Lemma~\ref{lem:ray} yield
\begin{equation}
 X_x=\tau',\qquad X_y=-3\tau,\qquad
 X_{g3a}=3t\tau\tau'.
 \label{eq:rcoordinates}
\end{equation}
Indeed the zero $g3a$-output is
$tX_{g3a}+(tX_x)(tX_y)=0$.  Thus the linear observable
$x+y+g3a$ equals $r(t)$ on this inverse.

Finally define
\begin{equation}
 P_{44}(A,B)=\ii\sum_{j=1}^{22}h_{4,j}(A+\ii B)B_j.
 \label{eq:P44}
\end{equation}
It is a homogeneous quartic with 538 expanded monomials.  At $(R_*,0)$ let
$\ell_{44}(A,B)=(x+y+g3a)(A+\ii B)$.  Lemmas
\ref{lem:symmetric-transfer} and~\ref{lem:coefficients} give
\begin{equation}
 D_{v_{44}}\bigl(\Delta_{44}^mP_{44}^{m+1}\bigr)(R_*,0)
 =2^m m!(m+1)!r_{m+1}\ne0
 \qquad(m\ge0).
 \label{eq:P44nonzero}
\end{equation}
This proves the quartic part of Theorem~\ref{thm:master} by a closed formula
at every order, strengthening the earlier collision-based conclusion of
nonvanishing at infinitely many unspecified orders.  This is the
homogeneous-quartic formulation equivalent to the Jacobian
conjecture~\cite{zhao-hnp}; the degree-eight companion instead addresses the
general homogeneous Vanishing Conjecture~\cite{zhao-general-vc}.

\section{Verification, scope, and archival policy}

The unified certificate contains the fully expanded 178-term and 608-term
potentials, the fixed inverse targets, the two coefficient formulas, the
fiber basis, and hashes of the precursor 14- and 44-variable sparse
certificates.  The primary exact checker verifies the block identities,
weighted conjugacy, exact fiber, $(Jg)^{13}g=0$, both inverse series,
symmetrization gradients, degrees, and term counts.  Its peak memory in the
reference run was below 100 MB.  A dependency-free standard-library checker
audits the exported sparse data and the first 1000 nonzero coefficients.
Finite coefficient checks are guards; the all-order proof is
Lemma~\ref{lem:coefficients}.

The formula $3+(4-2)+(6-2)+(7-2)=14$, the stable degree reductions, and the
symmetric reduction are classical.  The Discovery~06 precursor contains a
coefficient-Hankel proof of optimality only inside its stated constant-state
realization ansatz; that scoped side result is not needed here.  No global
smallest-dimension or smallest-term assertion is made.  Zihan Zhang had
already recorded existential Image and quartic Vanishing consequences of the
announced counterexample~\cite{zhang}; the contribution sought here is an
explicit fixed-dimensional, every-order realization.  At the audit cutoff,
no earlier public artifact was located with this exact 14-variable map, an
explicit $\SIC(n)$ witness for $n\le14$, or the every-order 44-variable
quartic formula.  That is provisional search evidence, not proof of worldwide
priority.

This paper is the canonical consequence paper.  The earlier 22/44-variable
note, the 21-variable SIC note, and the 14-variable construction note remain
public as timestamped technical precursors and certificate sources; they are
not presented as three additional publications.  The separate iterated
monodromy paper studies different mathematics and is not absorbed here.

\section*{AI assistance and responsibility}

The constructions, audits, programs, certificates, website, and manuscript
were developed with heavy assistance from ChatGPT 5.6 Sol under Alec
Kriebel's direction.  Alec Kriebel is the named human author and has requested
prominent disclosure that he is not qualified to validate the mathematics
independently.  Independent expert scrutiny is essential.

\begin{thebibliography}{99}
\bibitem{alpoge}
L. Alp\"oge, X post announcing the three-variable map, crediting Akhil Mathew
for posing the question and Claude Fable for producing the example,
20 July 2026;
\url{https://x.com/__alpoge__/status/2079028340955197566}.

\bibitem{bcw}
H. Bass, E. H. Connell, and D. Wright,
The Jacobian conjecture: reduction of degree and formal expansion of the
inverse, \emph{Bull. Amer. Math. Soc.} 7 (1982), 287--330.

\bibitem{campbell}
L. A. Campbell,
Unipotent Jacobian matrices and univalent maps,
\emph{Contemp. Math.} 264 (2000), 157--177;
\url{https://doi.org/10.1090/conm/264/04218},
\url{https://arxiv.org/abs/math/9907157}.

\bibitem{chamberland}
M. Chamberland and G. Meisters,
A mountain pass to the Jacobian conjecture,
\emph{Canad. Math. Bull.} 41 (1998), 442--451;
\url{https://doi.org/10.4153/CMB-1998-058-4}.

\bibitem{debondt}
M. de Bondt and A. van den Essen,
A reduction of the Jacobian conjecture to the symmetric case,
\emph{Proc. Amer. Math. Soc.} 133 (2005), 2201--2205;
\url{https://doi.org/10.1090/S0002-9939-05-07570-2}.

\bibitem{kulikov}
V. S. Kulikov,
The Jacobian conjecture and nilpotent maps,
\emph{J. Math. Sci.} 106 (2001), 3312--3319;
\url{https://doi.org/10.1023/A:1017907526453},
\url{https://arxiv.org/abs/math/9803143}.

\bibitem{liusun}
D. Liu and X. Sun,
Images of higher-order differential operators of polynomial algebras,
\emph{Bull. Aust. Math. Soc.} 96 (2017), 205--211;
\url{https://doi.org/10.1017/S0004972717000454}.

\bibitem{zhao-ag}
W. Zhao,
New proofs for the Abhyankar--Gurjar inversion formula and the equivalence of
the Jacobian conjecture and the vanishing conjecture,
\url{https://arxiv.org/abs/0907.3991}.

\bibitem{zhao-general-vc}
W. Zhao,
A vanishing conjecture on differential operators with constant coefficients,
\emph{Acta Math. Vietnam.} 32 (2007), 259--286;
\url{https://arxiv.org/abs/0704.1691}.

\bibitem{zhao-hnp}
W. Zhao,
Hessian nilpotent polynomials and the Jacobian conjecture,
\emph{Trans. Amer. Math. Soc.} 359 (2007), 249--274;
\url{https://doi.org/10.1090/S0002-9947-06-03898-0},
\url{https://arxiv.org/abs/math/0409534}.

\bibitem{zhao-image}
W. Zhao,
Images of commuting differential operators of order one with constant
leading coefficients,
\emph{J. Algebra} 324 (2010), 231--247;
\url{https://doi.org/10.1016/j.jalgebra.2010.04.022},
\url{https://arxiv.org/abs/0902.0210}.

\bibitem{zhang}
Z. Zhang,
Direct consequences of the three-dimensional counterexample to the Jacobian
conjecture, 20 July 2026;
\url{https://zzhang-iu.github.io/papers/direct-consequences-jacobian/}.
\end{thebibliography}

\end{document}
